Erwartet war (F-139), dass die Spitzen an den Verbinderenden sinken und $\sigma_{99}$ um 5--15~\% fällt. Das Gegenteil
trat ein: Mit \SI{50}{GPa} statt \SI{183}{GPa} steigt $\sigma_{99}$ im Gehäuse von \SI{308}{MPa} auf \SI{358}{MPa}
($+16$~\%), die Verformung von \SI{19}{mm} auf \SI{27}{mm}; nur die Verbinder selbst werden entlastet
(\SI{192}{\to}\SI{127}{MPa}). Der Grund: Die Verbinder sind hier der Entlastungsweg. Macht man sie insgesamt weich,
wandert die Last zurück in die Biegung der Gehäusewände. Die Idee bleibt richtig, muss aber \emph{örtlich} angewendet
werden: weiche Enden (wo die Spitze sitzt) und ein steifer Kern, der die Last weiter trägt. Dafür braucht das Modell
Verbinder mit abgestufter Steifigkeit -- der nächste Schritt.
